% PROVENANCE: independently audited by root 2026-11-19 14:09 UTC; incorporated into part_B.tex.
\paragraph*{Archived drafting provenance; authoritative version is in part B.}
\begin{saresult}{B.038}{Proved}{A derivation obstruction to tensor injectivity}
Let $R$ be a domain with fraction field $K$, let $V\subset K$ be
a discrete valuation overring of $R$, and let $D:K\to K$ be a
derivation satisfying $D(R)\subset V$. Suppose there are
$u_n\in R$ with $v(u_n)\to\infty$ and $D(u_n)=1$.
For $A=\operatorname{Int}(R)$ the map
\[
 \lambda:A\longrightarrow V,\qquad F\longmapsto F'(0)
\]
is well-defined. There is an $R$-balanced bilinear map
\[
 \beta:A\times A\longrightarrow K/V,
 \qquad \beta(F,G)=D(F'(0))G'(0)+V.
\]
If $P\in R[X]$, $c\in K$, $cP\in A$ and
$D(c)P'(0)^2\notin V$, then
\[
 (cP)\otimes P-P\otimes(cP)
\]
is a nonzero tensor with zero polynomial image.
\end{saresult}
\begin{proof}
Apply $D$ to $F(u_n)\in R$. The polynomial chain rule gives
\[
 D(F(u_n))=(DF)(u_n)+F'(u_n),
\]
where $DF$ means applying $D$ to the coefficients.
Its constant coefficient is $D(F(0))\in V$.
All nonconstant terms of $(DF)(u_n)$ have nonnegative valuation
for sufficiently large $n$, since their coefficients are fixed
and $v(u_n)\to\infty$. Hence $(DF)(u_n)\in V$ eventually.
The left side lies in $V$, so $F'(u_n)\in V$.
The nonconstant terms of $F'(u_n)-F'(0)$ likewise have
nonnegative valuation eventually. Therefore $F'(0)\in V$.

Additivity of $\beta$ is immediate. For $r\in R$,
\[
 D(rF'(0))G'(0)-rD(F'(0))G'(0)
   =D(r)F'(0)G'(0)\in V.
\]
Thus $\beta(rF,G)=r\beta(F,G)=\beta(F,rG)$, so it induces
an $R$-linear functional on $A\otimes_RA$ with target $K/V$.
On the displayed tensor its value is
\[
 \bigl(D(cP'(0))P'(0)-D(P'(0))cP'(0)\bigr)+V
   =D(c)P'(0)^2+V,
\]
which is nonzero by assumption. Its polynomial image is zero
by cancellation of the two equal products. This proves the
obstruction without any flatness assumption.
\end{proof}

\begin{saresult}{B.039}{Proved}{An explicit nonfinite-normalization counterexample}
Let $k=\mathbb F_2$, $A_0=k[t]_{(t)}$, $n_0=0$ and
$n_i=(i+1)!$ for $i\ge1$. In $k[[t]]$ set
\[
 z=\sum_{i\ge0}t^{n_i},\qquad
 z_i=\frac{z-\sum_{j<i}t^{n_j}}{t^{n_i}},\qquad
 u=t(z-1),\qquad s_i=(z_i-1)^2.
\]
Define
\[
 C=A_0[u,s_0,s_1,\ldots],\qquad
 M=(t,u)C,\qquad R=C_M.
\]
Then $R$ is a one-dimensional Noetherian local domain with
residue field $k$, fraction field $K=k(t,z)$ and normalization
a discrete valuation ring $V$ with parameter $t$ and residue
field $k$. Its normalization is not a finite $R$-module.
Nevertheless, unlike the finite-conductor one-branch cases,
$\operatorname{Int}(R)$ has a noninjective tensor map.
Explicitly, with $q(X)=X^2-X$ and $c=z-1$,
\[
 cq\in\operatorname{Int}(R),\qquad
 (cq)\otimes q-q\otimes(cq)\ne0,
\]
while its polynomial image is zero.
\end{saresult}
\begin{proof}
This is an explicit characteristic-two specialization of the
Akizuki--Reid construction. The known ring-theoretic conclusions
are stated in Reid, \emph{Akizuki's counterexample}, the two
theorems on page 2 and their proofs on pages 2--4; the opened
author LaTeX source is lines 49--212. We verify the hypotheses
and give the short characteristic-two arguments needed below,
including a direct nonfiniteness proof.

First $z$ is transcendental over $k(t)$. Otherwise let
$0\ne P(t,Z)\in k[t,Z]$ vanish at $z$, with $Z$-degree $d$
and coefficient $t$-degrees at most $h$. Its truncation
$z^{(N)}=\sum_{i\le N}t^{n_i}$ gives
\[
 v_t(P(t,z^{(N)}))\ge n_{N+1},\qquad
 \deg_t P(t,z^{(N)})\le h+d n_N.
\]
For all sufficiently large $N$, $n_{N+1}>h+dn_N$, so the
polynomial $P(t,z^{(N)})$ is zero. The distinct elements
$z^{(N)}\in k(t)$ would then give infinitely many roots of the
nonzero polynomial $P(t,Z)$ over that field, a contradiction.
The gap condition $n_i\ge2n_{i-1}+2$ of Reid's theorem also
holds, including $i=1,2$.

Here are direct checks of the local-ring assertions in this
specialization. Put $B=A_0[z_0,z_1,\ldots]$. The equations
\[
 z_i-1=t^{n_{i+1}-n_i}z_{i+1}
\]
make $tB$ the kernel of the residue map $B\to k$.
The local domain $V=B_{tB}$ has principal maximal ideal $tV$
and $\bigcap_n t^nV=0$, because it embeds in $k[[t]]$.
Every nonzero element therefore equals a power of $t$ times
a unit. Every nonzero ideal has an element of least such power,
so is principal; this proves that $V$ is a discrete valuation
ring, without presupposing Noetherianity of $B$.

In characteristic two,
\[
 s_i=t^{2(n_{i+1}-n_i)}(1+s_{i+1}),\qquad u^2=t^2s_0.
\]
Thus every $s_i$ lies in $tC$, and $C/M=k$.
Iterating the first identity shows that, modulo $t^NC$,
each of finitely many specified $s_i$ is represented by an
element of $A_0$. Using the second identity to reduce powers
of $u$ shows that $C/t^NC$ is generated over $A_0/(t^N)$ by
$1,u$. More precisely, with
$u_r=t(z-\sum_{j\le r}t^{n_j})=t^{n_r+1}(z_r-1)$,
every $f\in C$ can be written
\[
 f=X+Yu_r+t^NZ,\qquad X,Y\in A_0,\quad Z\in C,
\]
for any prescribed $r,N$, since $u-u_r\in A_0$.
The radical of $tC$ is $M$, as $C/tC$ is generated by the
square-zero image of $u$. Hence localization at $M$ does not
change $C/t^NC$, which is a finite $A_0/(t^N)$-module.

For $0\ne f\in M$, let $v_t(f)=a$, choose $N>a$ and then
$r$ with $n_r\ge N$. In the preceding expression $X$ has
valuation $a$. Divide by its unit factor in $A_0$ to write
\[
 f=t^a(1+t^{N-a}Z)+Yu_r.
\]
Squaring in characteristic two and using
$u_r^2=t^{2n_r+2}s_r$ gives
$f^2=t^{2a}w$ with $w\in C\setminus M$.
Thus $fR$ contains $t^{2a}$. The same conclusion holds for
any nonzero element of $R$ after clearing its unit denominator.
Every nonzero ideal consequently contains a power of $t$;
its quotient by that principal ideal is a submodule of a
finite $A_0/(t^N)$-module, and is finitely generated. This proves
that $R$ is Noetherian. Its only nonzero prime is $MR$, so its
dimension is one and its residue field is $k$.

Both $R$ and $V$ have fraction field $K$, since $z=1+u/t$.
Every $z_i$ is integral over $C$, as $z_i^2=1+s_i\in C$.
After first localizing $B$ by $C\setminus M$, all its maximal
ideals lie over $MR$ and contain $t$, so the only one is the
extension of $tB$. This localized ring is therefore $B_{tB}=V$
and is integral over $R$. As $V$ is integrally closed in $K$,
it is the normalization.

The derivation $D=t^{-1}\partial/\partial z$ of $K$ preserves
$R$: it kills $A_0$ and all the squares $s_i$, sends $u$ to $1$,
and extends to the localization by the quotient rule.
If $V$ were a finite $R$-module with generators $v_1,\ldots,v_h$,
then the valuations of $D(V)$ would be uniformly bounded below
by the minimum of $0,v(D(v_1)),\ldots,v(D(v_h))$.
Indeed $D(rv_j)=D(r)v_j+rD(v_j)$ and $D(r),r\in R\subset V$.
But $z_i\in V$ and
\[
 D(z_i)=t^{-n_i-1},
\]
whose valuations are unbounded below. This proves nonfiniteness
directly.

The elements $u_r\in R$ satisfy $D(u_r)=1$ and
$v_t(u_r)=n_{r+1}+1\to\infty$, so B.038 applies.
Also $cMR\subset R$, since $ct=u$ and $cu=t s_0$ both lie in
$R$. As $q(R)\subset MR$ by the residue field $k=\mathbb F_2$,
we have $cq\in\operatorname{Int}(R)$.
Finally $D(c)=1/t\notin V$ and $q'(0)=-1$.
The functional of B.038 takes the displayed tensor to
$1/t+V\ne0$, whereas its polynomial image is zero.
This proves the asserted counterexample with all its hypotheses.
\end{proof}


\begin{saresult}{B.040}{Proved}{Every finite-residue Akizuki--Reid family fails tensor uniqueness}
Let $A_0$ be any discrete valuation ring with uniformizer $t$ and
finite residue field of cardinality $q$. Choose units $a_i\in A_0$
and exponents $n_0=0$, $n_{i+1}\ge2n_i+2$, such that
$z=\sum_i a_it^{n_i}\in\widehat A_0$ is transcendental over
$\operatorname{Frac}(A_0)$. Define $z_i,C,M,R$ by Reid's construction:
\[
 z_i=\frac{z-\sum_{j<i}a_jt^{n_j}}{t^{n_i}},\qquad
 C=A_0[t(z-a_0),(z_i-a_i)^2:i\ge0],\quad
 M=(t,t(z-a_0))C,\quad R=C_M.
\]
Then $R$ is a one-dimensional Noetherian local domain with
nonfinite discrete valuation normalization $V$, as in Reid's
opened theorem. With $c=z-a_0$ and $P=X^q-X$, one has
$cP\in\operatorname{Int}(R)$ and
\[
 (cP)\otimes P-P\otimes(cP)\ne0,
\]
although its polynomial image is zero. This includes mixed
characteristic; it is not restricted to characteristic two.
\end{saresult}
\begin{proof}
The ring assertions, with exactly the hypotheses just stated,
are Reid's two theorems on page 2, with proofs on pages 2--4.
Set $D=t^{-1}\partial/\partial z$ on
$K=\operatorname{Frac}(A_0)(z)$. This derivation kills $A_0$
and sends $t(z-a_0)$ to $1$. On the remaining generators,
\[
 D((z_i-a_i)^2)=2t^{-n_i-1}(z_i-a_i).
\]
Its valuation is at least
$n_{i+1}-2n_i-1\ge1$; if $2=0$ the value is zero.
Thus $D(C)\subset V$. Denominators from $C\setminus M$ are
units of $V$, so the quotient rule gives $D(R)\subset V$.
The elements
\[
 u_i=t\left(z-\sum_{j\le i}a_jt^{n_j}\right)
      =t^{n_i+1}(z_i-a_i)\in R
\]
satisfy $D(u_i)=1$ and $v_t(u_i)=n_{i+1}+1\to\infty$.
Consequently B.038 applies with its weaker, explicitly stated
hypothesis $D(R)\subset V$.

Writing $u=t(z-a_0)$, one has $ct=u$ and
$cu=t(z-a_0)^2\in C$, so $cMR\subset R$.
The polynomial $P$ has values in $MR$ because the residue field
has $q$ elements. Hence $cP\in\operatorname{Int}(R)$.
Finally $D(c)=1/t$ and $P'(0)=-1$; B.038 takes the tensor to
$1/t+V\ne0$ and proves the assertion.
The same bounded-derivative argument as in B.039 reproves
nonfiniteness of $V/R$, since $D(z_i)=t^{-n_i-1}$ has
unboundedly negative valuation.
\end{proof}
